\section{Related Work}
\label{sec:related_work}

\subsection{Monocular Depth Estimation}
Early approaches to single-image depth estimation framed the task as continuous regression using convolutional neural networks~\cite{eigen2014depth}. Self-supervised methods emerged to learn depth from stereo pairs~\cite{godard2019digging} or monocular video sequences without ground-truth LiDAR. Ranftl et al.~\cite{ranftl2020towards} established large-scale multi-dataset training via affine-invariant loss functions ($\text{MiDaS}$). 

Recently, vision transformers (ViTs) have become the predominant paradigm for dense spatial prediction~\cite{ranftl2021vision,bhat2021adabins}. Foundation models such as Depth Anything~\cite{yang2024depth,yang2024depth2} and Metric3D/Metric3D v2~\cite{yin2023metric3d,yin2024metric3dv2} achieve powerful visual representations by pairing massive transformer backbones ($>300$\,M to $1.1$\,B parameters) with tens of millions of curated images and self-supervised DINOv2 visual features~\cite{oquab2023dinov2}. Concurrently, ZoeDepth~\cite{bhat2023zoedepth} and Depth Pro~\cite{bochkovskiy2024depthpro} explored metric scale transfer and high-resolution patch stitching. In parallel, feedforward geometric models such as MoGe~\cite{wang2024moge} explore multi-view consistent geometric representation. 

Alongside heavy foundation models, a distinct line of research addresses lightweight, edge-capable depth models: FastDepth~\cite{wofk2019fastdepth} introduced mobile-optimized encoder-decoder topologies, MonoViT~\cite{zhao2022monovit} combined self-attention with convolutional blocks, and Lite-Mono~\cite{zhang2023litemono} demonstrated effective self-supervised depth estimation in the $3.1$\,M to $8.5$\,M parameter envelope. Dioptra-DINO addresses the gap between these regimes: coupling a compact self-supervised DINOv2-Small representation with explicit camera-geometric inductive biases within an edge-feasible parameter budget ($27.5$\,M parameters, $105.05$\,MB weights).

\subsection{Camera-Aware Geometric Encodings}
Standard vision transformers treat tokens as discrete spatial coordinates $(x, y)$ using 2D sinusoidal or learned positional embeddings~\cite{vaswani2017attention,dosovitskiy2020image}. While effective for semantic recognition, 2D coordinate embeddings lack physical optical grounding: altering camera field-of-view (FOV) or principal point alters image contents without modifying the 2D coordinate grid.

To incorporate camera geometry, CamConvs~\cite{facil2019camconvs} appended normalized ray coordinates as auxiliary convolutional channels. Guizilini et al.~\cite{guizilini20203d} developed 3D convolutions for point cloud packing, while Metric3D~\cite{yin2023metric3d,yin2024metric3dv2} introduced virtual focal length canonicalization. UniDepth and UniDepthV2~\cite{piccinelli2024unidepth,piccinelli2025unidepth2} formulated camera-aware tokenization to predict pseudo-spherical metric representations across disparate camera models. While UniDepth targets multi-dataset foundation scale ($>300$\,M parameters), Dioptra-DINO develops an ultra-compact, continuous ray-unprojection formulation that operates directly within lightweight token modulation and attention bias computation.

\subsection{Lightweight Spatial Perception on Edge Hardware}
Deploying dense 3D visual perception on micro-aerial vehicles (MAVs), robots, and mobile devices requires strict adherence to memory and thermal budgets. While recent attention optimizations and LayerScale~\cite{touvron2021going} enhance training stability, peak memory utilization during high-resolution activation passes remains a key bottleneck on shared unified memory architectures (such as Apple Silicon and embedded NVIDIA Jetson platforms). Dioptra-DINO enforces a compact memory footprint ($27.5$\,M parameters, $105.05$\,MB FP32 weights, $52.5$\,MB in FP16), operating comfortably within a $<210$\,MB peak RAM envelope on embedded robotics platforms.
